\section{Normal geometry and admissible phases}\label{sec:preliminaries}

We begin with the geometric information used by both corrections. A nonzero
second fundamental form supplies enough directions for the small oscillations;
it need not satisfy any rank condition beyond being nonzero as a tensor.

Fix a smooth background metric and a finite collection of precompact coordinate
patches on $M$. All $C^m$ norms and unmarked tangent norms use these fixed data.
For a smooth immersion $F:M\to\R^4$, write
\[
 \gamma(F)=F^*\delta_4,
 \qquad N_F=(dF(TM))^\perp,
 \qquad B_F(V,W)=\operatorname{proj}_{N_F}D_V(dF(W)).
\]
The last expression is tensorial; in coordinates $B_{F,ij}$ is the normal
projection of $F_{ij}$. Let $K_\gamma$ denote the Gaussian curvature of
$\gamma=\gamma(F)$. Our sign convention for the Gauss equation is
\begin{equation}\label{eq:gauss}
 \langle B_F(V,V),B_F(W,W)\rangle-|B_F(V,W)|^2
 =K_\gamma\,|V\wedge W|_\gamma^2.
\end{equation}
In coordinates $x,y$, the right side is $K_\gamma\det\gamma$ for
$V=\partial_x,W=\partial_y$. The Hessian convention gives
$\Hess_\gamma x(\partial_y,\partial_y)=-\Gamma^x_{yy}$.
To check \eqref{eq:gauss}, choose Riemannian normal coordinates for
$\gamma$, orthonormal at the point under consideration.
There $\Gamma=0$, so $F_{xx},F_{xy},F_{yy}$ are normal vectors, and the
coordinate curvature formula is
\[
 K_\gamma=\partial_{xy}\gamma_{xy}
 -\tfrac12\bigl(\partial_{xx}\gamma_{yy}
                  +\partial_{yy}\gamma_{xx}\bigr).
\]
Substituting $\gamma_{ij}=\langle F_i,F_j\rangle$ cancels the third
derivatives and gives
$K_\gamma=\langle F_{xx},F_{yy}\rangle-|F_{xy}|^2$ at that point.
Tensoriality gives \eqref{eq:gauss} for arbitrary $V,W$.

A \emph{preferred normal} is a smooth unit section of $N_F$. This is less than a
normal frame: we will use complementary normal vectors only within coordinate
disks. None of the following local constructions requires an orientation of $M$.

\begin{definition}
A nonzero covector $\xi\in T_p^*M$ is \emph{good for $F$} if
$B_F(Y,Y)\ne0$ for a nonzero vector $Y\in\ker\xi$.
This condition is independent of the choice of $Y$ in that line.
\end{definition}

If $B_F(p)\ne0$, there are only finitely many bad projective covectors at $p$.
Indeed, take a normal component of $B_F(p)$ that is a nonzero scalar quadratic
form. Its zeros in the projective tangent line are finite, and a zero of the
vector-valued quadratic form must be one of these zeros.

We choose these good directions for the small oscillatory corrections.
They need not be the phases $x_j$ of the prescribed primitive additions in
\eqref{eq:intro-sum}: those phases will satisfy a transverse Hessian
condition, and their transverse second fundamental form may vanish in the
interior of the support disk.

Writing a positive metric as a sum of positive rank-one forms is a basic
part of Nash's construction \cite{Nash1954}; see
\cite[\S3.5.4, p.~32 and footnote~30]{Gromov2015} and
\cite[Lemmas~2.2--2.3]{Kallen1978} for related decompositions.
Here the individual phases and their relevant pairwise combinations must
also be good for the current second fundamental form.  The following
construction supplies those additional conditions.

\begin{lemma}[Positive phase decompositions]\label{lem:phase-data}
Let $F:M\to\R^4$ be a smooth immersion with $B_F(p)\ne0$ for every $p$, and let
$h_0$ be a smooth positive definite tensor. There are finitely many functions
$\phi_\nu$ defined on coordinate patches, smooth cutoffs $\psi_\nu$ supported
compactly in those patches, and smooth fiberwise linear forms $q_\nu$ on
symmetric tensors, such that
\begin{equation}\label{eq:phase-decomposition}
 H=\sum_\nu b_\nu(H)^2\,d\phi_\nu^2,
 \qquad b_\nu(H)=\psi_\nu\sqrt{q_\nu(H)}
\end{equation}
whenever $H$ is sufficiently close to $h_0$ in $C^0$. Every square root has a
strict positive lower bound before multiplication by the cutoff. On the
appropriate supports the differentials of $\phi_\nu$, $2\phi_\nu$, and
$\phi_\nu\pm\phi_\mu$ for $\nu\ne\mu$ are nonzero and good. These conclusions
persist for immersions in a sufficiently small $C^2$ neighborhood of $F$.
Each phase or phase combination has a smooth complementary coordinate
coframe on the support where it is used.
\end{lemma}

\begin{proof}
At a point choose an $h_0$-orthonormal coframe and three covectors with angles
$0,\pi/3,2\pi/3$, rotated by a common angle. Their squares form a basis of the
three-dimensional space of symmetric tensors, and $h_0$ is a positive linear
combination of them. A rotation avoids the finite bad set simultaneously for
all three directions. Extend the covectors as differentials of coordinate-linear
functions on a small patch. The basis, positivity, and goodness conditions
persist after shrinking the patch. The three coefficient forms in this basis
are smooth in position. A finite such cover and a partition
$\sum_\alpha\psi_\alpha^2=1$, with each cutoff supported in its patch, give
\eqref{eq:phase-decomposition}, using three indices over each patch.

It remains to control combinations from different indices. Multiply successive
phases by sufficiently rapidly increasing positive constants, and divide their
coefficient forms by the squares of those constants. For two intersecting
supports, the differential of either sum or difference, after normalization,
is arbitrarily close to that of the member with the larger weight. There are
only finitely many supports, and goodness has a positive angular margin on
their compact closures. Inductively chosen weights therefore make every
required combination good and nonzero. Doubles and signs change no kernel
line. A coordinate complementary to the dominant phase remains complementary
to the combination. All conditions are strict conditions on the first two
jets of the immersion; compactness gives their persistence.
\end{proof}

The small-oscillation argument is local on these coframes. If a combination is
not globally injective as a coordinate function, the dual coordinate vector
fields still make sense where its differential and the chosen complement are
independent. Subdivision gives actual coordinate neighborhoods when necessary.
Amplitudes are compactly supported there, so extension by zero is smooth.

For the exact correction, the phase data must also be controlled when the
input immersion varies with a small parameter. We record this quantitative
version separately. In the statement, a quantity is \emph{polynomially bounded}
if its absolute value is at most $Cz^{-A}$ for some finite $A\ge0$; the exponent
may depend on a specified finite derivative order, but not on $z$.

\begin{lemma}[Phase data for a varying family]\label{lem:polynomial-phase-data}
Fix a metric $h_*$ and suppose that $f_z$, $0<z<z_0$, are smooth
immersions satisfying uniform upper and positive lower metric bounds,
\[
 \|f_z\|_{C^j}=O(z^{-j})\quad(0\le j\le3),
 \qquad
 \inf_{p\in M}\max_{|Y|=1}|B_{f_z}(Y,Y)|\ge c>0.
\]
Then the phase data of Lemma~\ref{lem:phase-data}, with $h_0=h_*$, can be chosen
on cells of diameter comparable to $z^6$. The number of cells and all fixed-order
derivatives of their cutoff and phase data are polynomially bounded.
The good-direction estimates and metric nondegeneracy persist throughout the
$C^2$ ball of radius $z^4$ about $f_z$, for all sufficiently small $z$.

No regularity with respect to $z$ is required.
The construction uses a bounded number of phase weights, independent of the
number of cells. The inverse coframes and the reciprocals of the relevant
second-form lengths have polynomial bounds. Consequently every fixed finite
differential construction formed from these data by sums, products, derivatives,
metric inversion, normal projection, and division by the stated nonzero
quantities has polynomial bounds, provided its input has the required finite
polynomial derivative bounds. Its exponents depend only on the derivative
orders and on the input exponents, not on their implicit constants.
\end{lemma}

\begin{proof}
Metric bounds imply $df_z=O(1)$. Differentiating the formula for normal
projection gives
\[
 B_{f_z}=O(z^{-2}),\qquad DB_{f_z}=O(z^{-4}).
\]
For example, $D\operatorname{proj}_{N_{f_z}}=O(z^{-2})$ and multiplication by
$D^2f_z=O(z^{-2})$ gives the latter exponent; the $D^3f_z$ term is smaller.
The constants depend only on the fixed background data, metric bounds, and
specified family bounds.

At each cell center normalize the second fundamental form to have norm one.
The space of normalized normal-valued symmetric forms, in background
orthonormal frames, is compact. For each such form a rotated positive triple
avoids all its bad directions. Thus the maximum, over rotations, of the least
of the three second-form lengths has a strictly positive lower bound over
this compact set. Shrinking a uniform angular neighborhood preserves a fixed
fraction of this bound. This supplies a uniform angular margin and an estimate
\[
 |B_{f_z}(Y,Y)|\ge\kappa |B_{f_z}|
\]
for normalized kernel vectors in the selected directions, with $\kappa>0$
independent of the cell. Fixed coordinate-frame comparisons only alter
$\kappa$ by a uniform factor.

Use grids of mesh $z^6$ in the fixed charts, restricted to fixed compact
subpatches covering $M$. Smooth cutoffs equal to one on the covering cells and
supported in their fixed enlargements have derivatives $O(z^{-6j})$ of order
$j$. Divide them by the square root of their sum of squares. This gives a
partition by squares with the same type of bounds. A fixed lattice coloring,
combined with the finite chart index, gives a bounded number of colors such
that enlarged supports of the same color are disjoint. The number of cells is
$O(z^{-12})$.

On an enlarged cell the second form varies by $O(z^2)$; its nonzero lower bound
makes this variation small relative to its length. The center's phase triple
therefore stays good. The fixed smooth tensor $h_*$ varies by $O(z^6)$, so the
positive basis decomposition has a uniform positivity margin. Assign one
successive weight to each color and each place in the triple. Intersecting
supports have different such labels. Since the angular margins are uniform,
a fixed finite sequence of dominant weights makes every pairwise combination
good. Derivatives of the resulting coframes and their inverses are bounded at
every fixed order by the fixed charts and these weights. No weight depends on
the total number of cells.

If $\|G-f_z\|_{C^2}\le z^4$, projection onto the tangent and normal spaces changes
by $O(z^4)$. Hence
\[
 |B_G-B_{f_z}|\le C z^4(1+z^{-2})=O(z^2),
\]
so the direction and metric margins persist. Local unit normals used in mode
solves are obtained by normalizing a perpendicular product of two tangent
vectors and a nonzero vector $B_G(Y,Y)$. Their denominators are bounded below
by the metric and direction margins. The coefficient forms in
\eqref{eq:phase-decomposition} retain positive arguments for their square roots
on a fixed $C^0$ neighborhood of $h_*$. Product and chain rules now prove the
last assertion: finitely many inversions and differentiations multiply only
finitely many polynomial bounds. The polynomial number of summands can be
absorbed in another finite exponent.
\end{proof}
